\documentclass[11pt,a4paper]{article}
\usepackage[T1]{fontenc}
\usepackage{lmodern}
\usepackage[a4paper,margin=2.4cm]{geometry}
\usepackage{amsmath,amssymb}
\usepackage{parskip}
\usepackage{microtype}

\title{\textbf{The Essence of the Cellular Automaton}\\[2mm]
\large A self-directed emergent spiral on a pulsating wavefront}
\author{}
\date{}

\begin{document}
\maketitle
\vspace{-8mm}

\section{State per cell}
Each cell of the $L^3$ lattice ($L = 201$ in the reference build; $321$
for the full-size run) stores only:
\[
\text{cell} \;=\; \big(\; r^2,\;\; r,\;\; \texttt{active},\;\;
\texttt{spin},\;\; w \;\big),
\]
where $r^2$ is the squared Euclidean distance to the center
($\texttt{INF}$ = not yet visited), $r$ is the integer radius, two
bits mark membership --- \texttt{active} (on the pulsating shell) and
\texttt{spin} (on the spiral arm) --- and $w \in [\,0,\,9L-1\,]$ is a
fixed per-cell value, hashed once from the coordinates at startup,
that feeds the direction election of Sect.~\ref{sec:election}.

\section{Pulsating wavefront --- the ``sum of odds''}
The distance to the center is never computed by a square root or by
multiplication: it propagates through BFS. A unit step along axis $i$
changes $r^2$ by exactly the corresponding odd number, because
\begin{equation}
(n+1)^2 - n^2 \;=\; 2n + 1 .
\end{equation}
Hence, for a cell $x$ and its neighbor $x + e_i$, where $c_i$ denotes the
coordinate centered on axis $i$:
\begin{equation}
r^2(x + e_i) \;\leftarrow\; \min\!\big(r^2(x+e_i),\;\; r^2(x) + 2\,|c_i| + 1\big).
\end{equation}

On top of this field acts the pulse: a \emph{triangular} wave in $r^2$
space,
\begin{equation}
p(t) \in \big[\,0,\;\; 0.92\,R_{\max}^2\,\big],
\qquad R_{\max} = \tfrac{L}{2},
\end{equation}
which expands and contracts over time. A cell belongs to the living shell
exactly when
\begin{equation}
\texttt{active}(x,t) \;\iff\; \big|\, r^2(x) - p(t) \,\big| \;\le\; \tau ,
\end{equation}
with a small tolerance $\tau = 1$. The $0.92$ factor keeps the whole
dynamics strictly inside the lattice: no active cell ever exceeds
$r = \lfloor\sqrt{0.92}\,R_{\max}\rfloor$.

\section{The helical walker (3D Bresenham around an arbitrary axis)}

\subsection{Target geometry}
A cylinder of radius $R_{\mathrm{cyl}} \approx 0.182\,R_{\max}$ whose axis
is the line $\{\,P + sA\,\}$, with
\begin{equation}
|A| = \tfrac{L}{2}, \qquad |P| = R_{\mathrm{cyl}}, \qquad P \perp A .
\end{equation}
Because the center lies on that cylinder surface, the helix starts at
$r = 0$ and ends at $r = L/2$ (polar angle $\theta : 0 \to \pi$): one full
half-turn while climbing from the core to the outer shell.

\subsection{Radius to the axis without multiplication}
For $u = v - P$ (offset from the cylinder axis line), the squared distance
to the axis is
\begin{equation}
d^2 \;=\; \frac{|u \times A|^2}{|A|^2},
\end{equation}
so the walker compares $|u \times A|^2$ against the fixed target
$\mathrm{TGT} = R_{\mathrm{cyl}}^2 |A|^2$. Both the cross product
$c = u \times A$ and its squared norm $E = |c|^2$ are maintained
\emph{incrementally}: a lattice move $e_m$ changes $c$ by the constant
vector
\begin{equation}
D_m \;=\; e_m \times A,
\end{equation}
and changes $E$ by
\begin{equation}
G[m] \;=\; 2\,(c \cdot D_m) + |D_m|^2,
\end{equation}
whose own update after move $k$ uses the precomputed table
\begin{equation}
K[k][m] \;=\; 2\, D_k \cdot D_m .
\end{equation}
This is the exact generalization of the classic ``$2dx + 1$'' increment
from scalar coordinates to cross-product space.

\subsection{Orbital step}
Among the six lattice moves whose tangential projection is positive
($t = -c$ enforces a single sense of rotation), pick the one minimizing
the radial error
\begin{equation}
\big|\, E + G[m] - \mathrm{TGT} \,\big|,
\end{equation}
restricted to the tolerance band $|\cdot| \le \mathrm{TOL}$ (about one
cell); ties favor the largest angular progress. If no move stays inside
the band, fall back to minimum radial error.

\subsection{Climb step}
A multi-dimensional DDA along the axis: each accumulator obeys
$\mathit{dda}_i \mathrel{+}= |a_i|$; the largest one fires the climb move
and then subtracts $|a_x| + |a_y| + |a_z|$. Lattice moves are therefore
distributed exactly in proportion to the axis components.

\subsection{Interleaving}
A Bresenham accumulator interleaves orbital and climb steps with period
\begin{equation}
\mathit{orbit\_total} + \mathit{climb\_total},
\end{equation}
tracing the helix one cell per tick.

\section{The constraint that defines the automaton's character}
At runtime only additions, subtractions, shifts, and comparisons exist:
\emph{no multiplication, no division, no floating point, no trigonometry}.
Everything that depends on the chosen axis --- $D_m$, $K[k][m]$, $P$,
$\mathrm{TGT}$, $\mathrm{TOL}$, $\mathit{orbit\_total}$ --- is derived once
at initialization time, where multiplication is allowed. All geometric
intelligence lives in precomputed constants plus local comparison rules.

\section{Emergence}
\texttt{spin} draws the arm; \texttt{active} paints the living shell.
The spiral itself is never programmed: it emerges from the local rule
``minimize radial error while keeping the rotation sense,'' executed on
top of the pulsating bubble.

\section{Electing the direction: the ant-view tournament on $w$}
\label{sec:election}

The fused automaton is not told where to spiral. At every expansion
limit of the pulse --- the exact tick where $p(t)$ turns from ascending
to descending --- the automaton harvests its own next rotation axis out
of the $w$ field, through a tournament that no cell referees.

\subsection{Sowing}
Every \texttt{active} shell cell is seeded with the payload
\begin{equation}
\mathrm{payload}(x) \;=\; \bigl(\mathrm{score}(x) \ll 24\bigr)\;\big|\;\mathrm{code}(x),
\qquad
\mathrm{score} = H\bigl((w(x) + s) \bmod 9L,\; \mathrm{code}(x)\bigr),
\end{equation}
where $\mathrm{code}(x)$ is the cell's linear index (low bits) and $s$
is the global seed dial. Because the code occupies the low bits, the
payload order is a \emph{strict total order}: ties are impossible, so
the winner is unique and independent of any scan order.

\subsection{Diffusion --- and the flood trap}
Each pass lets every cell adopt the greatest payload among its six face
neighbors,
\begin{equation}
\mathrm{payload}(x) \;\leftarrow\; \max\bigl(\mathrm{payload}(x),\;
\mathrm{payload}(x \pm e_i)\bigr),
\end{equation}
alternating sweep directions; $2R + 8$ passes guarantee full coverage of
the torus. The subtlety: after enough passes the winning payload has been
\emph{copied} across a large region of the grid, down to the lowest
linear indices. An index scan for the maximum would therefore decode a
\emph{copy} --- possibly the corner cell $(0,0,0)$, yielding
$|m| = \sqrt{3}\,L/2$ and a nonsense axis. The implementation avoids the
trap structurally: the true winner --- the \emph{seeded} cell holding the
greatest payload --- is identified once, at sowing time, when payloads
are still unique. The diffusion then remains what it should be: pure
transport of the news, never a redefinition of the winner.

\subsection{From winner to axis}
The winner becomes the momentum vector: the raw displacement from the
lattice center to the winning cell,
\begin{equation}
\boldsymbol{m} \;=\; x_{\mathrm{win}} - c .
\end{equation}
Only its \emph{direction} carries meaning; the magnitude is discarded by
rescaling,
\begin{equation}
A \;=\; \frac{L/2}{|\boldsymbol{m}|}\,\boldsymbol{m},
\end{equation}
so the helix always walks on an axis of length exactly $L/2$. The axis is
frozen while a walk is alive; the next election (next expansion limit)
prepares a fresh direction for the next run. Nothing is prescribed: the
direction is harvested from the CA's own state at the exact instant the
wave reaches its limit.

\section{This edition: the ant-view broadcast}
An additional layer follows the same philosophy. Each cell holds one
monotonically growing value ($0$ = never reached). When the walker
computes a new spiral point, that cell receives the current tick as its
value. Then, every other tick (sweep direction alternating), every cell
inspects \emph{only} its six face neighbors and copies the greatest
neighbor value if it exceeds its own:
\begin{equation}
v(x) \;\leftarrow\; \max\big(v(x),\; v(x \pm e_i)\big).
\end{equation}
No cell ever senses beyond its neighborhood, yet the values diffuse as
waves that in finite time reach every cell of the $L^3$ grid --- global
coverage achieved purely by local moves, in lockstep with the CA tick.

Rendered as faint dust over a sparse sample, the \emph{hue} encodes
\emph{which expansion wave} delivered the value --- one distinct,
coherent colour per breathing cycle,
\begin{equation}
\mathrm{hue} \;=\; \Bigl\lfloor \frac{v}{T_{\mathrm{half}}} \Bigr\rfloor
\cdot \varphi \;\bmod\; 1,
\qquad \varphi = \tfrac{1+\sqrt5}{2},
\end{equation}
where $T_{\mathrm{half}}$ is the tick length of one expansion phase of
the pulse --- while the \emph{brightness} encodes how recently the value
arrived, glowing for a bounded youth window and then settling to a dim
floor.

\end{document}